% DEFINED:   sec:three-axes, sec:richer-dgp
% RESOLVED:  sec:inductive-biases, sec:parallel-axes, ch:cnn, ch:rnn,
%            ch:fixed-context-lm (backward into Part I and earlier Part II)
% FORWARD (prose-only): Chapter 7 (interpretability), Chapter 8 (RL).
%            Not drafted yet; stays prose.

\section{Inductive bias, three axes}\label{sec:three-axes}

Across the book we have walked the two axes of supervised inductive bias
named in \Cref{sec:inductive-biases}:

\begin{itemize}
\item \textbf{Architecture} (this chapter and the previous two,
  \Cref{ch:cnn} and \Cref{ch:rnn}): the structural prior the network
  commits to about \emph{how features compose}. The CNN commits to
  locality and translation equivariance; the RNN to time-translation
  equivariance and a bottleneck state; the transformer to a
  content-addressable lookup, repeated and composed through depth.
\item \textbf{Output head} (\Cref{sec:parallel-axes}): the prior the
  interpretation commits to about \emph{the distribution of the response}.
\end{itemize}

The pointer experiment forces a third dimension into the open:

\begin{itemize}
\item \textbf{Implementation realization}: even when the right inductive
  bias is structurally present in the architecture, the choice of
  initialization, positional encoding, supervision shape, and optimizer
  can keep it from materializing during training. At production scale,
  with the standard defaults, this dimension is mostly invisible. At
  small experimental scale it becomes the dominant factor, and it is
  where most of the debugging cycles in real transformer engineering go.
\end{itemize}

Every successful supervised network in modern practice is a triple: a
body that bakes in the right architectural prior for the data, a head
that bakes in the right likelihood prior for the response, and a training
recipe co-evolved with the body and head so the priors can actually be
discovered by gradient descent. The combinations multiply.

The transformer is the most flexible body in the book. It bakes in less
structure than a CNN or an RNN: it does not commit to spatial locality,
to temporal recurrence, or to a fixed context window. It commits instead
to a \emph{mechanism}, content-addressable retrieval, and lets the data
dictate which positions are retrieved when. The cost of that flexibility
is the $T \times T$ attention matrix, quadratic in the sequence length in
both compute and memory. The benefit is that one architecture serves any
task whose structure is ``attend to the right earlier position.''
Language modeling is the canonical case (the right position depends on
the syntax, the topic, the entity in play), and the same prior fits code
completion, reasoning over context, and retrieval-augmented generation.

The cost is \emph{also} the realization burden. The flexibility is
genuinely accessible only with a co-evolved recipe, and much of modern
transformer practice (warmup, learned or rotary positional encodings,
careful initialization, layer-norm placement, layer-wise learning-rate
decay, mixed precision, gradient clipping) is the accumulated body of
tricks that keeps the inductive bias usable. \Cref{sec:pe-scale-fix}
rediscovered one slice of that recipe by debugging.

The next chapter takes the model just trained and reverse-engineers its
circuit, asking whether the architectural argument of
\Cref{sec:one-layer-limit} is empirically vindicated and how the circuit
relates to induction heads and in-context learning. The model is small
enough to interpret fully, and the answer clarifies what the
content-addressable-lookup primitive actually is.

The closing chapter introduces reinforcement learning, the third
learning paradigm beyond supervised. The heads-as-bias frame still holds
there (a policy head is a softmax over actions; a value head is identity
plus mean squared error on returns), but the training signal stops being
a per-example label and becomes a scalar reward over trajectories. The
inductive-bias frame extends naturally: reward shaping, policy
architecture, and exploration strategy are all priors. The theoretical
optimum, AIXI, is Solomonoff induction (the compression-as-prediction
north star of \Cref{ch:fixed-context-lm}) joined to Bayesian decision
theory and reward maximization.

\section*{Appendix: a richer DGP that did not pan out}\label{sec:richer-dgp}

The first design for this chapter used a recursive bit-stream generator
with prefix-coded instructions (\texttt{examples/bit\_dgp.py}): a
continuous bit stream in which each instruction was either a literal bit
(prefix code \texttt{0} for value 0, \texttt{10} for value 1) or a
dereference (\texttt{11} followed by an Elias-gamma-encoded offset and the
dereference's own value codeword), with the offset counted in instruction
units. Multi-hop chains arose naturally whenever a dereference's target
was itself a dereference.

The generator is rich and pedagogically appealing: the model must learn
the prefix code, parse the stream into instructions, count instructions
for the gamma-encoded address, and perform the lookup, all from raw bits.
Small transformers (model width 64 to 128, one or two layers) did not
learn it within our training budget; even at 5000 iterations they barely
beat the marginal-prediction baseline. Parsing, counting, and addressing
in one small model appears to need more capacity or far longer training
than we had. The positional-encoding-scale problem of
\Cref{sec:pe-scale-fix} likely contributed, and we did not retry the
recursive generator with a learned encoding.

We kept the generator as reference because the design is interesting in
its own right: it is the smallest recursive task we know that exercises
stacked attention end to end with no task-specific scaffolding. Showing
that a sufficiently large or sufficiently well-tuned transformer can
learn it is left open.
